\chapter{Introducción}\label{cap:intro}

Aquí va la introducción del capítulo (motivación, objetivos, hipótesis,
preguntas de investigación, contribución y organización de la tesis).

\section{Motivación}\label{sec:intro-motivacion}
% - RI neuronal dominada por el inglés; el español desatendido
% - Rerankers efectivos pero costosos; la fusión de rangos como alternativa sin cómputo neuronal adicional
% - MessIRve habilita la evaluación nativa en español

\section{Objetivos}\label{sec:intro-objetivos}
% - Objetivo general y los seis objetivos específicos del protocolo, cada uno mapeado a su capítulo/experimento

\section{Hipótesis}\label{sec:intro-hipotesis}
% - Una hipótesis por pregunta de investigación, derivadas del protocolo

\section{Preguntas de investigación}\label{sec:intro-preguntas}
% - Las cuatro (P1–P4), mapeadas a los experimentos E1–E4

\section{Contribución}\label{sec:intro-contribucion}
% - Evidencia sistemática de fusión de rangos en español
% - Comparación de seis algoritmos de fusión
% - Análisis de complementariedad léxico-semántica
% - Sistema abierto y reproducible

\section{Organización de la tesis}\label{sec:intro-organizacion}
% - Un párrafo por capítulo
